\documentclass[11pt]{article}
\usepackage{amsmath,hyperref}
\title{Returning to contact before spending the cave relay}
\author{}\date{}
\begin{document}\maketitle

\paragraph{What carries over.}
The parent relay check holds both senders fixed. That is not a fair
null if the custodians can move before sending. This check allows
movement on the published cave graph, still using the weighted
visibility graph and the parent experiment's hard radio gates:
cost at most 10 and a declared 72 m range. Origin-to-label lookup,
unit-time movement, and every physical edge closure are modeled.
A hard-gate pass gives named one-slot service in this model; the
calculation does not measure packet delivery or passage clearance.

\paragraph{Costs and benchmark.}
Custodians start at 18 and 16, the fixed receiver occupies 39, and
an installer at 13 has one optional relay. A mobile actor can move
one edge, wait or send during a slot. Senders cannot move while
sending; the receiver admits one report per slot. A send at
slot $t$ settles at its end, before the activation effective in
slot $t+1$. Installing a relay takes one slot at its chosen site,
may be delayed, and is irreversible. Each custodian retains
its report until settlement. The egress obligation applies to
both custodians and the installer. There is no other resource
or motion conflict, report transfer, optional inspection or
unverified fallback in this benchmark.

With fixed installed site and cut, feasibility is equivalent
to choosing two distinct send slots and corresponding contact
vertices within the outbound and remaining egress budgets, with
a feasible installer path. Necessity follows by reading those
cuts from the run; sufficiency follows from the available shortest
walks and waits around the send slots. An independent time-expanded
recurrence tracks positions before and after each admissible send.
It agrees on 2,400 finite deadline queries: four cut patterns, five
placement or motion scopes, and 120 triples per arm. The most
permissive placement arm allows any of the 158 tile origins that
the installer can reach in time, not just the two prior sites.

\paragraph{What actually needs the relay.}
For completion cuts $(1,2,9)$ (old report, new report, all
egress), no relay is needed with the alternate mouth open:
$A$ sends from 18 in slot 0, $B$ moves to 42 and sends in slot 1.
That schedule still fits when only $2$--$3$ is closed.
If only $13$--$42$ is closed, $B$ can make the slot-1 send
from 42 without a relay, but has no time to egress by 9.
A no-relay return-to-contact schedule reaches egress 9 if $B$'s
report deadline is 5 instead; $B$ walks to 2 before sending.
With relay 13 effective in slot 1, $A$ can settle before the
change and $B$ can move to 6, send in slot 1, and egress by 9.
Fixing both sender positions rejects that feasible case.
Permitting all placements rather than only 13 does not add
any feasible deadline triple in the finite sweep. Closing both
gates traps $B$ in the movement graph; radio does not supply an exit.

\paragraph{Recipient and timing of the gate observation.}
Compare an open alternate mouth to a closed one with the
short gate open in both. They remain in that condition for
the horizon. A qualified read at 13 tells the installer which
world holds; this is an assumed local input, not a detector
event. The two local histories at $B$ at the start of slot 0
are identical without an extra signal. Its first motion has
to be chosen then. When closed, only position 6 at the end
of slot 0 supports the three deadlines, and an effective
relay is necessary. When open, only 42 works without a
relay; with relay 13, both 6 and 42 work. The installer has
no time to travel and make a different installation effective
in slot 1. Thus a blind common choice is 6, which spends
the relay even in the open world. A closure bit available to
$B$ before slot 0 instead selects 42 without spending in the
open world, and 6 with spending in the closed world.
A perfect report at slot 1 cannot restore the option. For
zero-error randomization the same obstruction holds on each
seed that satisfies both histories; their finite intersection
has probability one when each succeeds with probability one.
The three Lean statements check the action-table step after
the motion calculation.

An equally informed conventional scheduler can use the
same decision. The result identifies an evidence recipient
and cutoff, not a verified way to sense, send or act on that
bit in the simulator. It calibrates matching for I2; no independent
acceptance mission or new operating envelope is claimed.
\end{document}
